% 08_oq1_meissner_zeeman.tex — Section 8: OQ-1 Meissner-Zeeman Calculation
% Phase 2 full-body migration — Task #210, Cycle NVB, 2026-03-21
% Source: open_questions/oq1_meissner_zeeman_calculation.tex (Task #145, 2026-03-20; authoritative)
% Law 16: quantitative claims carry source locator comments.
% Status: POPULATED (Phase 2 full body) — full body extracted from oq1_meissner_zeeman_calculation.tex
%
% NOTE: \thebibliography block removed — master document handles bibliography via references/references.bib.
%       All \cite{} keys retained. \needref{} macros retained.
%       \jacobdecision macro not present in this source — no substitution needed.
%
% ARCHIVE: Previous section content (summary) archived as 08_oq1_meissner_zeeman_phase2_task210_archive_2026-03-21.tex

\section{OQ-1: Meissner--Zeeman Detection Feasibility}
\label{sec:oq1}

\subsection{Physical Setup and Assumptions}
\label{sec:oq1_setup}

The D0 demonstration places an NV diamond directly on (or adjacent to)
an NbN thin film on sapphire substrate.
A small permanent magnet provides a bias field $B_\text{bias}$ to lift the
$\pm m_s$ degeneracy and allow Zeeman shift detection.
When the sample cools through $T_c$, the NbN enters the Meissner state and
the local field at the NV changes.

\textbf{Key parameters (all from references or Jacob's data):}
\begin{itemize}
  \item NV gyromagnetic ratio: $\gamma_{NV} = 28.024$\,GHz/T \cite{Doherty2013}
        % LOCATOR: Doherty2013 Table 1 (NV Hamiltonian parameters)
  \item Zero-field splitting at 12\,K: $D = 2877.5$\,MHz (This work, Jacob's data)
  \item ODMR linewidth: $\Delta\nu_\text{line} = 10$\,MHz (This work)
  \item ODMR CW sensitivity: $\eta = 14\,\mu\text{T}/\sqrt{\text{Hz}}$ (This work)
        % LOCATOR: Jacob's measured sensitivity; see sensitivity_calculator.py and lab notebook entry 2026-03 (calibration run to be archived). Barry2020 provides the CW sensitivity formalism (Eq. 6 therein); the 14 µT/√Hz value is Jacob's own measured result, not from Barry2020.
  \item NbN London penetration depth: $\lambda_L \approx 400$\,nm
        \cite{Kamlapure2010}
        % [NVB 2026-03-19] Resolved: Kamlapure et al. APL 96, 072509 (2010).
        % PDF on disk: open_questions/references/Kamlapure_2010_NbN_PenetrationDepth_APL.pdf
        % Note: Kamlapure reports lambda_L(0) increasing with decreasing thickness
        % in epitaxial NbN on MgO; thickest (50nm) films with Tc~15.87K have
        % lambda in the 300-500nm range. Value 400nm is within reported range.
        % See also Chockalingam2008 (PRB 77, 214503) for NbN film properties.
        % SUBSTRATE CAVEAT: Kamlapure2010 NbN on silicon substrate; our substrate is sapphire — λ_L value used as order-of-magnitude estimate
  \item NbN film thickness: $d_\text{SC} = 200$\,nm (design target for D0)
  \item NbN $T_c \approx 15.7$\,K (This work, to be measured)
\end{itemize}

\subsection{Zeeman Shift Formula}
\label{sec:oq1_zeeman}

The NV center Hamiltonian (in the secular approximation) for a magnetic
field $\mathbf{B}$ projected onto the NV symmetry axis gives transition
frequencies \cite{Barry2020,Doherty2013}:
\begin{equation}
  f_\pm = D \pm \gamma_{NV}\, B_\text{NV}
  \label{eq:oq1_zeeman}
\end{equation}
where $B_\text{NV}$ is the projection of the local field on the NV axis.
The Zeeman splitting is $\Delta f = f_+ - f_- = 2\,\gamma_{NV}\,B_\text{NV}$.

When the SC transitions at $T_c$, the local field changes by $\Delta B_\text{NV}$,
giving an ODMR frequency shift:
\begin{equation}
  \Delta\nu_\text{signal} = 2\,\gamma_{NV}\,\Delta B_\text{NV}
  \label{eq:oq1_shift}
\end{equation}

\textbf{Detection threshold} for $\text{SNR} = 1$ in integration time $\tau$:
\begin{equation}
  \Delta B_\text{min}(\tau) = \frac{\eta}{\sqrt{\tau}}
  = \frac{14\,\mu\text{T}/\sqrt{\text{Hz}}}{\sqrt{\tau}}
  \label{eq:oq1_threshold}
\end{equation}

This is far more sensitive than the naive threshold $\Delta\nu > \Delta\nu_\text{line}$
(which would require $\Delta B > 0.357$\,mT) --- with 1\,s integration, the true
detection limit is $14\,\mu$T.

\subsection{Meissner Screening: Pearl Thin-Film Model}
\label{sec:oq1_pearl}

\subsubsection{Pearl Length}

For a thin SC film ($d_\text{SC} < 2\lambda_L$), the relevant screening length
is the Pearl length \cite{Pearl1964}% Pearl1964 p.65: thin-film vortex current sheet; Pearl length Λ = 2λ²/d introduced
:
\begin{equation}
  \Lambda = \frac{2\lambda_L^2}{d_\text{SC}}
  \label{eq:oq1_pearl}
\end{equation}
For NbN with $\lambda_L = 400$\,nm and $d_\text{SC} = 200$\,nm:
\begin{equation}
  \Lambda = \frac{2 \times (400\,\text{nm})^2}{200\,\text{nm}} = 1600\,\text{nm} = 1.6\,\mu\text{m}
\end{equation}

\subsubsection{Field Above the Thin SC Film}

For an infinite thin SC film in a uniform perpendicular field $B_\text{bias}$,
using the Pearl model, the $z$-component of the field at height $z$ above
the film center is \cite{Kogan2007}:
% MODEL CAVEAT: Kogan (2007) Pearl-length model valid when λ_L >> d (thin-film limit); confirm this holds for our 200nm NbN film
% GEOMETRY CAVEAT: Pearl (1964) assumes infinite thin film; our NbN-on-sapphire is finite geometry — Pearl model used as best available analytical framework
% NOTE: Kogan2007 Eq. 6 gives the Fourier-space vortex potential; the Meissner uniform-field screening formula below follows from the z >> Λ limit applied to the Pearl screening length — Pearl1964/Kogan2007 primarily derive vortex field distributions; this formula applies the same Pearl length Λ to the Meissner (vortex-free) screening regime.
\begin{equation}
  B_z(z) \approx B_\text{bias}\left(1 - \frac{\Lambda}{2z}\right) \quad (z \gg \Lambda)
  \label{eq:oq1_pearl_field}
\end{equation}

The field \emph{change} at the NV when the SC transitions ($T\to T_c^-$):
\begin{equation}
  \boxed{
  \Delta B_z(z) = -B_\text{bias}\,\frac{\Lambda}{2z}
  = -B_\text{bias}\,\frac{\lambda_L^2}{z\,d_\text{SC}}
  }
  \label{eq:oq1_deltaB}
\end{equation}

\textbf{Physical interpretation:} The thin SC film partially screens $B_\text{bias}$
from the region below the film (the SC interior).
At height $z \gg \Lambda$, the screening is weak (film is thin relative to distance).
At $z \sim \Lambda$, the full Meissner exclusion is felt.

\subsubsection{Numerical Estimates}

For $B_\text{bias} = 1$\,mT and $\Lambda = 1.6\,\mu$m:

\begin{table}[h!]
\centering
\caption{Expected Meissner field change $|\Delta B_z|$ vs.\ NV height $z$
  (from Eq.~\ref{eq:oq1_deltaB}; $B_\text{bias} = 1$\,mT, $\Lambda = 1.6\,\mu$m,
  $\eta = 14\,\mu\text{T}/\sqrt{\text{Hz}}$).
  Integration time $\tau = (\eta/|\Delta B_z|)^2$.}
\label{tab:oq1_feasibility}
\small
\begin{tabular}{@{}lllll@{}}
\toprule
$z$ (NV height) & $|\Delta B_z|$ & $\Delta\nu_\text{signal}$ & $\tau$ for SNR$=$1 & Feasible? \\
\midrule
$1\,\mu$m   & $0.80$\,mT  & $44.8$\,MHz & $<1$\,ms & \checkmark\checkmark \\
$10\,\mu$m  & $80\,\mu$T  & $4.5$\,MHz  & $31$\,ms & \checkmark\checkmark \\
$100\,\mu$m & $8\,\mu$T   & $0.45$\,MHz & $3.1$\,s & \checkmark \\
$500\,\mu$m & $1.6\,\mu$T & $0.09$\,MHz & $77$\,s  & \checkmark (marginal) \\
\bottomrule
\end{tabular}
\end{table}

\subsection{Required Bias Field}
\label{sec:oq1_bias}

For detection in time $\tau$ with NV at height $z$:
\begin{equation}
  B_\text{bias} > \frac{\eta}{\sqrt{\tau}} \cdot \frac{2z}{\Lambda}
  = \frac{\eta}{\sqrt{\tau}} \cdot \frac{z\,d_\text{SC}}{\lambda_L^2}
\end{equation}

For $z = 10\,\mu$m, $\tau = 1$\,s:
\begin{equation}
  B_\text{bias} > 14\,\mu\text{T} \times \frac{2 \times 10\,\mu\text{m}}{1.6\,\mu\text{m}}
  = 14\,\mu\text{T} \times 12.5 = 175\,\mu\text{T} \approx 1.75\,\text{G}
\end{equation}
% CORRECTED 2026-03-20 (Task #145): Previous value 87.5 µT used z/Λ = 6.25 instead of 2z/Λ = 12.5 — factor-of-2 error. Correct minimum bias is 175 µT. D0 recommendation (B_bias ≤ 10 mT) remains unchanged and well above this threshold.

A bias field of 0.1--1\,mT is easily achievable with a small NdFeB magnet
at 2--10\,cm standoff.

\subsection{Geometry Considerations}
\label{sec:oq1_geometry}

The calculation above assumes an \emph{infinite} SC sheet and \emph{perpendicular}
applied field. Practical considerations:

\begin{enumerate}
  \item \textbf{Film size:} For a finite disk of radius $R$ with $R \gg z$,
    Eq.~\eqref{eq:oq1_deltaB} is valid at the disk center.
    For a 5\,mm$\times$5\,mm NbN patch with $z = 10\,\mu$m: $R/z = 250 \gg 1$ \checkmark

  \item \textbf{Field orientation:} Eq.~\eqref{eq:oq1_pearl_field} applies to the
    perpendicular ($z$-direction) field component. An in-plane field creates
    a different (and generally weaker) response for a thin film. \emph{Bias magnet
    should be oriented with significant normal component.}

  \item \textbf{NV axis angle:} The Zeeman shift in Eq.~\eqref{eq:oq1_zeeman} uses
    the projection of $\Delta\mathbf{B}$ onto the NV symmetry axis.
    For NV $\langle 111 \rangle$ axis at angle $\theta$ from $z$:
    $\Delta B_\text{NV} = \Delta B_z \cos\theta$. For $\theta = 54.7^\circ$ (worst
    case for diamond (100) surface): $\cos\theta = 0.577$ --- factor of $\sim$2 reduction.

  \item \textbf{Vortex entry (above $H_{c1}$):} If $B_\text{bias} > \mu_0 H_{c1}$ for NbN,
    flux vortices enter and the Meissner state is only partial. NbN $H_{c1}$ is low
    ($\sim$20\,mT for bulk NbN; thin films may differ) \cite{Ito2019}.
    % SOURCE CAVEAT: Ito2019 H_c1 from multilayer SRF context; used as best available proxy for thin NbN film
    % LOCATOR: Ito2019 Fig. 6 shows effective H_c1 of NbN-SiO2-Nb multilayers; bulk NbN µ0H_c1 ≈ 20 mT drawn from general NbN literature consistent with Ito2019 context
    % [NVB 2026-03-19] Resolved: Ito et al. SRF2019 TUP078.
    % PDF on disk: open_questions/references/Ito_2019_NbN_Hc1_SRF2019.pdf
    % IMPORTANT NOTE: Literature consistently reports mu0*Hc1 ~ 20 mT (200 Oe)
    % for bulk NbN. The original estimate of ~1 mT in this document was incorrect by ~20x. Updated to ~20 mT.
    % This changes the D0 constraint: bias field up to ~20 mT is safe for
    % pure Meissner state, not just <<1 mT.
    For $B_\text{bias} \lesssim \mu_0 H_{c1} \approx 20$\,mT: pure Meissner state,
    full screening. \emph{Use $B_\text{bias} \lesssim 10$\,mT for D0.}
\end{enumerate}

\subsection{Conclusion and OQ-1 Status}
\label{sec:oq1_conclusion}

\textbf{The D0 Meissner--Zeeman demonstration is feasible} with Jacob's
current system:

\begin{itemize}
  \item With $B_\text{bias} = 0.1$\,mT and NV at $z = 10\,\mu$m:
    $|\Delta B_z| = 8\,\mu$T, detectable in $\tau = 3$\,s with $\eta = 14\,\mu\text{T}/\sqrt{\text{Hz}}$

  \item If bulk diamond (NV at $z = 100\,\mu$m): need $B_\text{bias} \gtrsim 1.0$\,mT and
    $\tau = 3$\,s --- still feasible with standard NdFeB magnet
    % CORRECTED 2026-03-20 (Task #145): Previous value 0.9 mT contained factor-of-2 error in B_bias > (η/√τ)(2z/Λ). Correct: (14/√3)×(2×100/1.6) ≈ 1.0 mT.

  \item \textbf{Critical geometry constraint:} Bias magnet must be oriented
    to provide $B \perp$ SC film ($B_z$ component). Use NV ensemble (all 4
    $\langle 111 \rangle$ orientations) for optimal coupling.

  \item \textbf{NbN parameter citations resolved (2026-03-19):}
    (a) NbN London depth $\lambda_L \approx 400$\,nm cited from \cite{Kamlapure2010};
    (b) NbN $H_{c1} \approx 20$\,mT cited from \cite{Ito2019}.
    Note: $\mu_0 H_{c1} \approx 20$\,mT for bulk NbN (not 1\,mT as originally estimated).
    The D0 bias field of $B_\text{bias} \lesssim 10$\,mT is safely below $H_{c1}$.
    \textbf{Needref status (updated 2026-03-19, Task \#98):}
    (c) Pearl 1964 (Eq.~\ref{eq:oq1_pearl}) --- RESOLVED: \cite{Pearl1964} [PDF in hand:
    \texttt{literature/Pearl1964\_superconducting\_films\_fluxoids.pdf}];
    % Pearl1964 p.65: introduces thin-film vortex fluxoid model; Pearl length Λ = 2λ²/d
    (d) Pearl model field formula (Eq.~\ref{eq:oq1_pearl_field}) --- RESOLVED: \cite{Kogan2007}
    [PDF in hand: Kogan, PRB 75, 064514 (2007); arXiv:cond-mat/0611187].
    Kogan explicitly introduces $\Lambda = 2\lambda^2/d$ and the Pearl vortex field model.
\end{itemize}

\textbf{Estimated Zeeman shift at $T_c$ crossing:}
\begin{equation}
  \Delta\nu_\text{Zeeman} \approx 2\,\gamma_{NV}\,B_\text{bias}\,\frac{\Lambda}{2z}
  = 2\,\gamma_{NV}\,B_\text{bias}\,\frac{\lambda_L^2}{z\,d_\text{SC}}
  \approx 0.45\,\text{MHz}
  \quad \text{(for $z=10\,\mu$m, $B_\text{bias} = 0.1$\,mT)}
\end{equation}
% CORRECTED 2026-03-20 (Task #145): Previous value 0.9 MHz used Λ/z instead of Λ/(2z) in numerical substitution — factor-of-2 error. Correct value: 2 × 28.024 MHz/mT × 0.1 mT × (0.4 µm)²/(10 µm × 0.2 µm) = 2 × 28.024 × 0.1 × 0.08 = 0.448 MHz ≈ 0.45 MHz. Consistent with Table 1, z=100 µm, B_bias=1 mT entry (same z×B_bias product gives same Δν).

\textbf{OQ-1 Status: RESOLVED.} The Meissner--Zeeman signal is detectable with Jacob's existing
apparatus at the D0 operating point. The only remaining blocker is the NbN film itself
(Section~\ref{sec:nbn-sourcing}).
